%\documentclass{article}
%\usepackage{verbatim}
%
%\begin{document}
\topic{Kecepatan Metode Gauss}
Kita menggunakan Metode Gauss untuk menyelesaikan sistem linear dalam buku ini
karena metode ini mudah dipahami, mudah dibuktikan memberikan jawaban yang
benar, dan cepat.
Metode ini cepat dalam arti bahwa, dalam semua perhitungan manual yang kita
perlukan, kita memperoleh jawabannya hanya dalam beberapa langkah yang memakan
waktu beberapa menit.
Namun,
ilmuwan dan insinyur yang dalam praktiknya menyelesaikan sistem linear harus
memiliki metode yang cukup cepat untuk sistem besar, dengan 1.000 persamaan,
10.000 persamaan, atau bahkan 100.000 persamaan.
Sistem-sistem ini diselesaikan dengan komputer, sehingga kecepatan mesin
membantu, tetapi kecepatan metode yang digunakan tetap menjadi pertimbangan
utama dan kadang-kadang menjadi faktor yang membatasi masalah mana yang dapat
diselesaikan.

Kecepatan suatu algoritma biasanya diukur dengan melihat bagaimana waktu yang
diperlukan untuk menyelesaikan masalah bertambah seiring bertambahnya ukuran
himpunan data masukan.
Dengan kata lain, berapa banyak waktu tambahan yang diperlukan algoritma jika
ukuran data masukan kita perbesar dengan faktor sepuluh, misalnya dari sistem
dengan 1.000 persamaan menjadi sistem dengan 10.000 persamaan, atau dari 10.000
menjadi 100.000?
Apakah waktu yang diperlukan menjadi sepuluh kali,
seratus kali, atau seribu kali lebih lama?
Apakah waktu yang diperlukan algoritma sebanding dengan ukuran himpunan data,
dengan kuadrat ukuran tersebut, dengan pangkat tiga ukuran tersebut, dan
seterusnya?

Berikut adalah penggalan kode Metode Gauss yang diimplementasikan dalam bahasa
komputer FORTRAN.
Koefisien sistem linear disimpan dalam larik $\nbyn{N}$ \textit{A}, sedangkan
konstantanya disimpan dalam larik $\nbym{N}{1}$ \textit{B}.
Untuk setiap \textit{ROW} antara $1$ dan $N$, program ini telah menemukan entri
utama $A(ROW,COL)$.
Selanjutnya, program akan menggabungkan baris-baris.
\begin{equation*}
  -\text{\textit{LEINV}}\cdot \rho_{\text{\textit{ROW}}}
      +\rho_{\text{\textit{i}}}
\end{equation*}
(Penggalan kode ini hanya untuk ilustrasi dan belum lengkap.
Sebagai contoh, lihat topik selanjutnya tentang Akurasi Metode Gauss.
Meskipun demikian, penggalan ini memadai untuk tujuan kita karena
analisis versi yang lengkap, termasuk semua pengujian
dan subkasus, lebih rumit tetapi pada dasarnya memberikan hasil yang sama.)
\begin{center}{\small
\begin{verbatim}
       DO 10 I=ROW+1, N
               LEINV=A(I,COL)/A(ROW,COL)
               DO 20 J=COL, N
                    A(I,J)=A(I,J)-LEINV*A(ROW,J)
                  20 CONTINUE
               B(I)=B(I)-LEINV*B(ROW)
          10 CONTINUE
\end{verbatim} 
}\end{center}

Perulangan terluar (tidak ditampilkan) melintasi $N-1$ baris.
Untuk setiap baris tersebut,
perulangan yang ditampilkan menjalankan operasi aritmetika pada entri-entri
dalam \textit{A} yang berada di bawah dan di sebelah kanan entri utama
(dan juga pada entri-entri dalam \textit{B}, tetapi untuk menyederhanakan
analisis, operasi tersebut tidak akan kita hitung---lihat \nearbyexercise{}).
Kita akan mengasumsikan bahwa entri utama ditemukan di posisi yang lazim,
yaitu
$\text{\textit{COL}}=\text{\textit{ROW}}$
(seperti di atas, analisis kasus umum lebih rumit tetapi pada dasarnya
memberikan hasil yang sama).
Ini berarti terdapat $(N-\text{\textit{ROW}})^2$ entri yang harus dikenai
operasi aritmetika.
Secara rata-rata, \textit{ROW} akan bernilai $N/2$.
Dengan demikian, kita memperkirakan perulangan bersarang di atas akan
dijalankan kira-kira $(N/2)^2$ kali,
yakni dalam waktu yang sebanding dengan kuadrat jumlah persamaan.
Dengan memperhitungkan perulangan terluar yang tidak ditampilkan, kita
memperoleh perkiraan bahwa waktu eksekusi algoritma sebanding dengan pangkat
tiga jumlah persamaan.

Algoritma yang berjalan dalam waktu yang berbanding lurus dengan ukuran
himpunan data tergolong cepat, algoritma yang berjalan dalam waktu yang
sebanding dengan kuadrat ukuran himpunan data lebih lambat tetapi biasanya
masih cukup layak digunakan, dan algoritma yang berjalan dalam waktu yang
sebanding dengan pangkat tiga ukuran himpunan data pun masih memiliki
kecepatan yang wajar.

Perkiraan kecepatan seperti ini merupakan cara yang baik untuk memahami
seberapa cepat atau lambat suatu algoritma diperkirakan berjalan secara
rata-rata.
Namun, pada beberapa sistem khusus, kode Metode Gauss di atas berjalan sangat
cepat. Karena itu, mungkin ada ciri-ciri suatu masalah yang menjadikannya
sangat cocok untuk jenis penyelesaian ini.

Dalam praktiknya, kode yang terdapat dalam sistem aljabar komputer atau paket
standar menerapkan suatu varian Metode Gauss yang disebut faktorisasi
segitiga.
Untuk menyatakan metode ini diperlukan bahasa aljabar matriks, yang baru akan
kita pelajari pada Bab Tiga.
Meskipun demikian, secara konseptual kode di atas cukup dekat dengan kode yang
biasanya digunakan dalam penerapan.

Ada beberapa percepatan teoretis dalam penyelesaian sistem linear.
Algoritma selain Metode Gauss telah diciptakan yang memerlukan waktu yang
sebanding bukan dengan pangkat tiga ukuran himpunan data, melainkan dengan
pangkat (kira-kira) $2.7$
(bidang ini masih aktif diteliti, sehingga eksponen tersebut mungkin akan
sedikit menurun seiring waktu).
Namun, peningkatan teoretis ini belum digunakan secara luas,
sebagian karena metode-metode baru tersebut memerlukan himpunan data yang
cukup besar sebelum dapat mengungguli Metode Gauss
(meskipun akan mengungguli Metode Gauss pada himpunan yang sangat besar,
terdapat biaya awal tertentu yang membuat metode-metode tersebut tidak lebih
cepat pada sistem yang, sejauh ini, telah diselesaikan dalam praktik).



\begin{exercises}
  \item Sistem komputer memungkinkan pembangkitan bilangan acak
    (tentu saja, bilangan-bilangan ini hanya acak semu karena dihasilkan
    oleh suatu algoritma, tetapi urutan bilangan yang diperoleh lolos
    sejumlah uji statistik yang wajar sehingga tampak acak).
    \begin{exparts}
      \partsitem Isi sebuah larik $\nbyn{5}$ dengan bilangan acak (misalnya,
        dalam rentang $[0..1)$).
        Terapkan Metode Gauss untuk menentukan apakah larik itu singular.
        Ulangi percobaan tersebut sepuluh kali.
        Apakah matriks singular sering atau jarang ditemukan (dalam pengertian
        ini)?
      \partsitem Ukur waktu yang dibutuhkan komputer untuk menyelesaikan
        sepuluh larik $\nbyn{5}$ berisi bilangan acak.
        Hitung waktu rata-ratanya.
        (Perhatikan bahwa beberapa sistem dapat diketahui singular dengan
        cukup cepat, misalnya jika baris pertama sama dengan baris kedua.
        Dengan mempertimbangkan bagian pertama, apakah Anda memperkirakan
        bahwa sistem singular berperan besar dalam rata-rata tersebut?)
      \partsitem Ulangi butir sebelumnya untuk larik $\nbyn{15}$.
      \partsitem Ulangi butir sebelumnya untuk larik $\nbyn{25}$.
      \partsitem Ulangi butir sebelumnya untuk larik $\nbyn{35}$.
      \partsitem Buat grafik ukuran masukan terhadap waktu rata-rata.
    \end{exparts}
  \item Larik $\nbyn{10}$ seperti apa yang dapat Anda rancang agar sistem
    komputer Anda memerlukan waktu paling lama untuk mereduksinya?
    Bagaimana dengan yang memerlukan waktu paling singkat?
  \item Lengkapi program FORTRAN untuk membuat implementasi langsung Metode
    Gauss.
    Bandingkan kecepatan kode Anda dengan kode yang digunakan dalam sistem
    aljabar komputer.
    Manakah yang lebih cepat?
    (Kebanyakan sistem aljabar komputer akan menerapkan beberapa teknik
    aljabar matriks yang akan kita pelajari nanti, pada Bab Tiga.)
  \item Perluas penggalan kode agar menangani kasus ketika larik \textit{B}
    memiliki lebih dari satu kolom.
    Dengan demikian, lebih dari satu sistem dapat diselesaikan sekaligus
    (semuanya dengan matriks koefisien \textit{A} yang sama).
  \item Spesifikasi bahasa FORTRAN mengharuskan larik disimpan ``berdasarkan
    kolom'', yaitu, seluruh kolom pertama disimpan secara kontigu, kemudian
    kolom kedua, dan seterusnya.
    Apakah penggalan kode yang diberikan memanfaatkan hal ini,
    atau dapatkah kode itu ditulis ulang agar lebih cepat (dengan memanfaatkan
    fakta bahwa pengambilan data oleh komputer berlangsung lebih cepat dari
    lokasi-lokasi yang kontigu)?
  \item Perkirakan waktu eksekusi reduksi Gauss-Jordan.
    Uji perkiraan Anda dengan mengimplementasikan reduksi Gauss-Jordan
    dalam suatu bahasa komputer dan menjalankannya pada matriks $\nbyn{5}$,
    $\nbyn{15}$, dan $\nbyn{25}$ dengan entri acak.
\end{exercises}
\endinput
